\section{Related Work}
\label{sec:related}

FAccT research has repeatedly shown that ``getting the metric right'' is insufficient when systems are embedded in institutions with asymmetric power, weak recourse, and opaque accountability.
We position \compa-\ecl-\mas at the intersection of (i) contestability and procedural justice for algorithmic systems, (ii) principle-to-practice governance, and (iii) model-level fairness tooling and trade-offs.

\begin{table*}[t]
\centering\small
\caption{Roles of related approaches; this conceptual positioning is not an exhaustive capability audit of all implementations.}
\label{tab:positioning}
\begin{tabularx}{\textwidth}{lXX}
\toprule
Approach & Main role & Relationship to this work \\
\midrule
Ethical principles & State institutional commitments & Motivate requirements, without fixing our schema. \\
Fairness toolkits & Measure or mitigate group disparities & Complement empirical ECL checks. \\
Impact assessments & Document impacts and responsibilities & Complement policy rationale and review procedures. \\
Contestability by design & Keep systems open to intervention & Motivates preserved case evidence and bounded replay. \\
COMPASS revision & Fixed-score policy audit and artifacts & Demonstrates explicit rules, reconstructed states and export gating. \\
\bottomrule
\end{tabularx}
\end{table*}


\subsection{Contestability, Due Process, and Institutional Accountability}
A central FAccT concern is whether affected people can \emph{understand, challenge, and obtain remedy} when algorithmic systems shape life chances.
Citron and Pasquale argue that scoring systems can erode due process by obscuring decision logic and shifting power toward institutions \cite{CitronPasquale2014ScoredSociety}.
Subsequent work emphasizes that accountability is not a property of an algorithm alone, but of the socio-technical system: what is documented, what is disclosed, who can audit, and what remedies exist \cite{Wieringa2020AccountFor,Metcalf2021AIA,Selbst2019Abstraction}.
In HCI and socio-technical design, contestability-by-design frameworks articulate practices that keep systems responsive to human intervention throughout the lifecycle \cite{Alfrink2023ContestableAI}.

Our contribution aligns with this line: we treat \emph{audit artifacts} (policy manifests, constraint reports, decision traces) as first-class outputs required for contestation and institutional oversight.

\subsection{From Principle-Level Ethics to Operational Governance}
Major AI governance frameworks---NIST AI RMF, OECD AI Principles, UNESCO, and EU Trustworthy AI guidelines---converge on high-level desiderata (fairness, transparency, accountability, robustness) but remain implementation-agnostic \cite{NISTAIRMF2023,OECD2019AI,UNESCO2021AI,EUHLEG2019Trustworthy}.
Scholars have highlighted the ``last-mile'' gap between ethical principles and operational mechanisms \cite{Dignum2018EthicsSpecialIssue,Mittelstadt2019Principles}.
We build on this insight, but focus on a specific public-sector domain where legitimacy and recourse are central: educational decision support.

\subsection{Fairness Tooling, Trade-offs, and Preference-Based Selection}
Model-level toolkits such as Fairlearn and Aequitas provide practical metrics, mitigation strategies, and reporting to support fairness evaluation \cite{Weerts2023Fairlearn,Saleiro2018Aequitas}.
These metric-oriented roles are complementary to our decision-evidence contract; we have not exhaustively audited every extension of these toolkits.
Trade-offs between utility, fairness, and other constraints are unavoidable; preference-based selection over Pareto-feasible candidates has been explored in multiple forms, including ``fairness without harm'' style decision rules and policy-oriented selection \cite{Ustun2019FairnessWithoutHarm}.
Our system-level contribution is to embed such choices into \emph{versioned policy} and to couple selection with a procedural safeguard (abstention to human review) and an artifact trail.

\subsection{Algorithmic Systems in Education}
Educational prediction systems (early warning, risk scoring, intervention targeting) are widespread, but educational data often contains measurement and historical biases that complicate fairness claims \cite{BakerHawn2022AlgorithmicBiasEducation}.
Our setting therefore emphasizes governance practices (documentation, auditability, and recourse) in addition to model metrics.


\paragraph{Abstention, selective prediction, and distribution-free uncertainty.}
Our triage setting uses abstention (``reject option'') to mark uncertain cases as requiring human review. This idea has a long history in statistical decision theory (e.g., Chow's rule) and has been revisited in modern \emph{selective classification} and distribution-free uncertainty quantification (e.g., SelectiveNet; conformal prediction) \cite{Chow1970Reject,Geifman2019SelectiveNet,Angelopoulos2021Conformal}. These lines of work motivate two design choices in our system: (i) treating high abstention as a \emph{governance signal} rather than a silent knob, and (ii) recording the abstention policy in a replayable artifact so that institutions can audit when and why automation requested human review.


\subsection{Provenance, reconstruction and integrity mechanisms}
Provenance and version tracking are established capabilities. PROV-O models entities, activities, agents and derivation relations \cite{Lebo2013PROVO}. ModelDB records models, pipelines, parameters and metrics \cite{Vartak2016ModelDB}; Ormenisan et al. use change capture and data versions for pipeline reproducibility and debugging \cite{Ormenisan2020TimeTravel}. Our score-only replay has a narrower scope than reconstructing inference from versioned features.

Cryptographic log guarantees also predate this work. RFC 9162 specifies signed transparency-log structures and inclusion/consistency proofs for certificates \cite{Laurie2021RFC9162}. Our unsigned batch roster and byte digests do not supply those guarantees. In relation to sociotechnical contestability \cite{Alfrink2023ContestableAI}, the present engine supports scrutiny but does not enact an institutional appeal. Appendix~\ref{app:mechanism-positioning} compares these specific mechanisms without treating unexamined implementations as lacking features. The contribution is a bounded, executable decision-policy contract and its evaluated distinctions, not a new provenance standard, log-security primitive or proof of first integration.

The contract in Section~\ref{sec:export-contract} makes the acceptance dependencies explicit and testable. It does not introduce a new provenance primitive: its evaluated claim is that reproducing available decisions can pass while complete evidence verification fails, and that verified evidence can still fail a declared export constraint. Table~\ref{tab:mechanism-positioning} situates the constituent mechanisms; Table~\ref{tab:fault-contracts} tests these separations in our implementation.
